\documentclass[12pt]{extarticle}

% ============================================================
% PACKAGES
% ============================================================

% Encodage et langue
\usepackage[french]{babel}

% Mise en page
\usepackage[a4paper, margin=2cm]{geometry}

\usepackage[francais,bloc,completemulti]{automultiplechoice}    

% Mathématiques
\usepackage{amsmath}
\usepackage{amssymb}

% Dessins
\usepackage{tikz}

% Listes
\usepackage{enumitem}

% Boîtes
\usepackage[most]{tcolorbox}

% Images
\usepackage{graphicx}

% Mise en page avancée
\usepackage{fancyhdr}
\usepackage{multicol}

% Tableaux
\usepackage{booktabs}
\usepackage{tabularx}
\usepackage{colortbl}

\usepackage[normalem]{ulem}

\usepackage{pgfplots}
\pgfplotsset{compat=1.18}
\usepackage{pgf-pie}


\usepackage{ifthen}

% ============================================================
% COULEURS
% ============================================================

% ------------------------------------------------------------
% Couleurs principales
% ------------------------------------------------------------

\definecolor{Bleu}{RGB}{40, 100, 180}
\definecolor{Vert}{RGB}{50, 150, 90}
\definecolor{Orange}{RGB}{230, 140, 40}

\definecolor{BleuClair}{RGB}{235, 245, 255}
\definecolor{VertClair}{RGB}{235, 250, 240}
\definecolor{OrangeClair}{RGB}{255, 245, 225}

% ------------------------------------------------------------
% Couleurs complémentaires
% ------------------------------------------------------------

% Rappels de cours — violet
\definecolor{Violet}{RGB}{125, 85, 170}
\definecolor{VioletClair}{RGB}{245, 238, 252}

% Théorèmes — rouge
\definecolor{Rouge}{RGB}{190, 55, 55}
\definecolor{RougeClair}{RGB}{255, 238, 238}

% Lemme — indigo
\definecolor{Indigo}{RGB}{75, 85, 155}
\definecolor{IndigoClair}{RGB}{240, 241, 252}

% Points méthode — turquoise
\definecolor{Turquoise}{RGB}{35, 145, 150}
\definecolor{TurquoiseClair}{RGB}{235, 250, 250}

% Erreurs / attention — jaune
\definecolor{Jaune}{RGB}{210, 165, 35}
\definecolor{JauneClair}{RGB}{255, 250, 225}

% ------------------------------------------------------------
% Neutres
% ------------------------------------------------------------

\definecolor{GrisClair}{RGB}{245, 245, 245}
\definecolor{Gris}{RGB}{100, 100, 100}


% ============================================================
% COMMANDES PERSONNALISÉES
% ============================================================

\newcommand{\bb}[1]{\textbf{#1}}
\newcommand{\ii}[1]{\textit{#1}}

\newcommand{\cm}{\,\mathrm{cm}}
\newcommand{\m}{\,\mathrm{m}}

\newcommand{\reponse}[1]{\boxed{#1}}

% ============================================================
% TABLEAUX
% ============================================================

\newcommand{\ligneentete}{%
    \rowcolor{Bleu}%
    \color{white}\bfseries
}

\newcommand{\entete}[1]{%
    \cellcolor{Bleu}\color{white}\bb{#1}%
}


\NewDocumentCommand{\hl}{O{JauneClair}m}{%
  \tcbox[
    on line,
    colback=#1,
    colframe=#1,
    boxrule=0pt,
    arc=4pt,
    left=0pt,
    right=0pt,
    top=1pt,
    bottom=1pt
  ]{#2}%
}

\NewDocumentCommand{\hlb}{O{JauneClair}m}{%
  \tcbox[
    on line,
    colback=#1,
    colframe=#1,
    boxrule=0pt,
    arc=4pt,
    left=0pt,
    right=0pt,
    top=1pt,
    bottom=1pt
    ]{\bb{#2}}%
}
\newlength{\answerwidth}


\newcommand{\answer}[2][VertClair]{%
    \ifmmode
        \settowidth{\answerwidth}{$\displaystyle #2$}%
        \ifanswers
            {\setlength{\fboxsep}{0pt}%
             \colorbox{#1}{$\displaystyle #2$}}%
        \else
            \makebox[0pt][l]{%
                \makebox[\answerwidth][l]{%
                    \leaders\hbox to 0.5em{\hss.\hss}\hfill
                }%
            }%
            \phantom{#2}%
        \fi
    \else
        \settowidth{\answerwidth}{#2}%
        \ifanswers
            \hl[#1]{#2}%
        \else
            \makebox[\answerwidth][l]{%
                \leaders\hbox to 0.5em{\hss.\hss}\hfill
            }%
        \fi
    \fi
}


\newcommand{\graphanswer}[1]{%
    \ifanswers
        #1%
    \else
        0%
    \fi
}

\newcommand{\exo}[1]{%
    \includegraphics[
        width=\linewidth,
        keepaspectratio
    ]{#1}%
}

\newcommand{\nodata}{}
\newcommand{\correctanswer}[1]{
    \explain{
            \newline\vspace{10pt}
            #1
    }
}

\newlength{\mylen}
\newtcolorbox{exercicebox}[1][\unskip]{
    colback=white,
    colframe=Bleu,
    boxrule=1pt,
    arc=3mm,
    title=\bb{Exercice #1}\AMCnumero{1},
    left=2mm,
    right=2mm,
    top=2mm,
    bottom=0mm
}

\newcounter{theorem}
\newtcolorbox{theoremebox}[1][\unskip]{
    colback=RougeClair,
    colframe=Rouge,
    boxrule=1pt,
    arc=3mm,
    title=\bb{Théorème \stepcounter{theorem}\thetheorem}\hspace{5pt}#1,
    left=5mm,
    right=5mm,
    top=4mm,
    bottom=4mm
}

\newcounter{probleme}
\newtcolorbox{probbox}[1][\unskip]{
    colback=RougeClair,
    colframe=Rouge,
    boxrule=1pt,
    arc=3mm,
    title=\bb{Problème \stepcounter{probleme}\theprobleme #1}\AMCnumero{1},
    left=5mm,
    right=5mm,
    top=2mm,
    bottom=2mm
}

\newenvironment{tableau}[1]
{%
    \renewcommand{\arraystretch}{1.3}%
    \begin{center}%
    \begin{tabularx}{\linewidth}{#1}%
    \toprule
}
{%
    \bottomrule
    \end{tabularx}%
    \end{center}%
    \renewcommand{\arraystretch}{1}%
}



% Set page layout after the config files
\geometry{
  paper=a4paper,
  hmargin=1cm,
  headheight=1cm,
  headsep=.3cm,
  footskip=0cm,
  top=1.5cm,
  left=1.3cm,
  right=1.3cm,
  bottom=1.5cm
}

%\geometry{hmargin=1cm,headheight=1cm,headsep=.3cm,footskip=0cm,top=1.5cm,bottom=1.5cm}

\def\AMCbeginQuestion#1#2{\textbf{(#1)\hspace{2mm}}}
\def\AMCotextReserved{\emph{(Reservé)}}
\AMCopenOpts{boxframerule=0pt,boxframerulecol=red,framerule=0pt}
\newdimen\AMCpostOquest\AMCpostOquest=-7mm
\AMCformHSpace=.3em

\begin{document}
\baremeDefautS{b=1,m=0}
\exemplaire{24}{    

% ---- Header ------
\fbox{
    \begin{minipage}{.43\linewidth}
        \namefield{Nom: }
        Prénom :\\
        Date:
    \end{minipage}
}
\hfill
\begin{minipage}{.48\linewidth}
    \AMCcodeGridInt[vertical=false]{etu}{2}
\end{minipage}
% ---- Fin Header ------


\begin{exercicebox}[\it{(Nombres relatifs)}]

    \begin{question}{sommeneg}
        \begin{minipage}{0.22\textwidth}
        $(-10)+(-2)$ vaut...
        \end{minipage}
        \hfill
        \begin{minipage}{0.69\textwidth}
        \begin{reponseshoriz}
            \mauvaise{8}
            \bonne{-12}
            \mauvaise{+20}
            \mauvaise{12}
        \end{reponseshoriz}
        \end{minipage}
    \end{question}


    \begin{question}{divposneg}
        $(-6)\div(-2)$ vaut…
        \begin{reponseshoriz}
            \mauvaise{-3}
            \mauvaise{-12}
            \bonne{+3}
            \mauvaise{-6}
        \end{reponseshoriz}
    \end{question}

    \begin{question}{subsub}
        $(-18)-(-8)$ vaut…
        \begin{reponseshoriz}
            \mauvaise{26}
            \bonne{-10}
            \mauvaise{-14}
            \mauvaise{-26}
        \end{reponseshoriz}
    \end{question}

    \begin{question}{prioneg}
        $9\times[\hspace{2pt}(-3)\times(+2)+5]$ vaut…
        \begin{reponseshoriz}
            \bonne{-9}
            \mauvaise{99}
            \mauvaise{59}
            \mauvaise{-81}
        \end{reponseshoriz}
    \end{question}


\end{exercicebox}

\begin{exercicebox}[\it{(Carré et racine carrée)}]

    \begin{question}{racine1}
        \begin{minipage}{0.22\textwidth}
        La racine carrée de deux est égale à :
        \end{minipage}
        \hfill
        \begin{minipage}{0.69\textwidth}
        \begin{reponseshoriz}
            \bonne{$\sqrt{2}$}
            \mauvaise{1.41}
            \mauvaise{4}
            \mauvaise{$2^2$}
        \end{reponseshoriz}
        \end{minipage}
    \end{question}

    \begin{question}{carré}
        \begin{minipage}{0.30\textwidth}
        Le carré de 12 est égale à :
        \end{minipage}
        \hfill
        \begin{minipage}{0.65\textwidth}
        \begin{reponseshoriz}
            \bonne{144}
            \mauvaise{124}
            \mauvaise{24}
            \mauvaise{3.5}
        \end{reponseshoriz}
        \end{minipage}
    \end{question}

    \begin{question}{carrépetit}
        \begin{minipage}{0.30\textwidth}
        Le carré de 1.5 est égale à :
        \end{minipage}
        \hfill
        \begin{minipage}{0.65\textwidth}
        \begin{reponseshoriz}
            \bonne{2.25}
            \mauvaise{3}
            \mauvaise{2.5}
            \mauvaise{2}
        \end{reponseshoriz}
        \end{minipage}
    \end{question}

\end{exercicebox}

\begin{exercicebox}[\it{(Pythagore)}]


    \begin{question}{qcmegapythagore2}    
        ABC est rectangle en C. L'égalité de Pythagore s'écrit :
        \begin{reponses}
            \mauvaise{$AB^2=BC+AC$}
            \mauvaise{$AB^2=AC+CB$}
            \mauvaise{$AC^2=BC^2+AB^2$}
            \bonne{$AB^2=AC^2+BC^2$}
        \end{reponses}
    \end{question}

    \begin{question}{qcmpythahypo}
        Les côtés de l'angle droit d'un triangle rectangle mesurent 3m et 4m.
        L'hypoténuse mesure :
        \begin{reponseshoriz}
            \mauvaise{3.5m}
            \mauvaise{7m}
            \mauvaise{5.5m}
            \bonne{5m}
        \end{reponseshoriz}
    \end{question}


    \begin{question}{qcmreciproque}    
        La réciproque du théorème de Pythagore sert à ...
        \begin{reponses}
            \mauvaise{montrer qu'un triangle n'est pas rectangle.}
            \mauvaise{calculer l'aire d'un triangle rectangle.}
            \mauvaise{calculer la longueur de l'hypoténuse dans un triangle
            rectangle.}
            \bonne{montrer qu'un triangle est rectangle}
        \end{reponses}
    \end{question}

\end{exercicebox}


}   
\end{document}
